% =====================================================================
\section{Цифровой вход: кнопки и прерывания}
% =====================================================================
\lessonintro
  {научиться читать состояние кнопки, понять роль подтягивающих резисторов, побороть дребезг, познакомиться с прерываниями}
  {урок 1 (выходы, \code{millis()}), урок 2 (вывод в Serial для проверки)}
  {кнопка на \pin{5} включает и выключает лампу}
  {тактовая кнопка}

\subsection{Проблема «висящего» входа}

Вход GPIO имеет очень высокое сопротивление. Если к нему подключить только кнопку на землю, то в отпущенном
состоянии вывод никуда не подключён и ловит наводки --- читается случайно \code{0} или \code{1}.
Решение --- \textbf{подтягивающий резистор}, задающий уровень по умолчанию.

\begin{figure}[H]
\centering
\begin{circuitikz}
  \draw (0,3) node[vcc]{3,3 В} to[R,l_=$R_\text{pu}$ \SI{45}{k\ohm}] (0,1) coordinate(a);
  \draw (a) to[short,-o] (1.5,1) node[right]{\pin{5}};
  \draw (a) to[push button, l_=S1] (0,-1) node[ground]{};
  \node[align=center,font=\sffamily\small] at (0.6,-2.3) {Pull-up (INPUT\_PULLUP)\\ отпущена = HIGH, нажата = LOW};
  \begin{scope}[xshift=7cm]
  \draw (0,3) node[vcc]{3,3 В} to[push button, l_=S1] (0,1) coordinate(b);
  \draw (b) to[short,-o] (1.5,1) node[right]{\pin{5}};
  \draw (b) to[R,l_=$R_\text{pd}$ \SI{45}{k\ohm}] (0,-1) node[ground]{};
  \node[align=center,font=\sffamily\small] at (0.6,-2.3) {Pull-down (INPUT\_PULLDOWN)\\ отпущена = LOW, нажата = HIGH};
  \end{scope}
\end{circuitikz}
\caption{Подтяжка к питанию и к земле. У ESP32-S3 оба резистора встроены (${\approx}\SI{45}{k\ohm}$).
В проекте используем pull-up: кнопка замыкает \pin{5} на землю}
\end{figure}

\begin{tip}
Встроенная подтяжка слабая. Для длинных проводов или в условиях помех добавьте внешний резистор
\SIrange{4.7}{10}{k\ohm}.
\end{tip}

\subsection{Опрос кнопки}

\codecap{Состояние кнопки в Serial Monitor}
\codeinput{l3_button_read}

Если просто написать «при \code{LOW} переключить лампу», лампа будет мигать всё время, пока кнопка
удерживается, и срабатывать по нескольку раз на одно нажатие. Нам нужно ловить именно
\emph{момент нажатия} --- и при этом учесть дребезг.

\subsection{Дребезг контактов}

Механические контакты при замыкании несколько раз «подпрыгивают». Микроконтроллер, опрашивающий вход
миллионы раз в секунду, видит это как серию нажатий.

\begin{figure}[H]
\centering
\begin{tikzpicture}[x=1cm,y=1cm]
  \draw[->] (0,0) -- (12.3,0) node[right]{$t$};
  \draw[->] (0,0) -- (0,2) node[above]{$U$};
  \draw (0.05,1.4) -- (-0.05,1.4) node[left,font=\scriptsize]{HIGH};
  \draw (0.05,0) -- (-0.05,0) node[left,font=\scriptsize]{LOW};
  \draw[very thick,accent] (0,1.4) -- (2,1.4) -- (2,0) -- (2.2,0) -- (2.2,1.2) -- (2.35,1.2) -- (2.35,0)
    -- (2.6,0) -- (2.6,0.9) -- (2.7,0.9) -- (2.7,0) -- (2.95,0) -- (2.95,0.5) -- (3.0,0.5) -- (3.0,0) -- (8,0)
    -- (8,1.4) -- (8.12,1.4) -- (8.12,0.3) -- (8.3,0.3) -- (8.3,1.4) -- (8.45,1.4) -- (8.45,0.8)-- (8.55,0.8) -- (8.55,1.4) -- (12,1.4);
  \draw[decorate,decoration={brace,amplitude=4pt}] (2,1.6) -- (3.0,1.6) node[midway,above=4pt,font=\scriptsize,align=center]{дребезг\\ 1--20 мс};
  \draw[decorate,decoration={brace,amplitude=4pt}] (8,1.6) -- (8.55,1.6) node[midway,above=4pt,font=\scriptsize]{дребезг};
  \node[font=\scriptsize] at (5.5,0.35) {кнопка нажата};
  \begin{scope}[yshift=-2.4cm]
  \draw[->] (0,0) -- (12.3,0) node[right]{$t$};
  \draw[->] (0,0) -- (0,2);
  \node[font=\scriptsize,anchor=south west] at (0.2,1.5) {после фильтра};
  \draw[very thick,okgreen] (0,1.4) -- (4.0,1.4) -- (4.0,0) -- (9.55,0) -- (9.55,1.4) -- (12,1.4);
  \draw[dashed,gray] (3.0,0) -- (3.0,4.0);
  \draw[dashed,gray] (8.55,0) -- (8.55,4.0);
  \draw[<->,esp] (3.0,0.7) -- node[above,font=\scriptsize]{50 мс} (4.0,0.7);
  \draw[<->,esp] (8.55,0.7) -- node[above,font=\scriptsize]{50 мс} (9.55,0.7);
  \node[font=\scriptsize,text=okgreen] at (4.0,-0.35) {событие «нажатие»};
  \end{scope}
\end{tikzpicture}
\caption{Дребезг контактов (сверху) и сигнал после программного фильтра (снизу)}
\end{figure}

Программный фильтр принимает новый уровень, только если он стабилен дольше \SI{50}{ms}. Оформим его
в функцию, которая возвращает \code{true} ровно один раз на каждое нажатие:

\codecap{Функция buttonPressed(): нажатие с подавлением дребезга}
\codeinput{l3_debounce}

\subsection{Шаг проекта 3: кнопка-выключатель}

\progress{3}

Подключаем кнопку между \pin{5} и GND. Функции \code{setLamp()}, \code{handleSerial()} и
\code{heartbeat()} остаются из шага 2, \code{buttonPressed()} --- из раздела выше. Меняются только
\code{setup()} и \code{loop()}:

\codecap{Умная лампа, шаг 3 (изменения относительно шага 2)}
\codeinput{lamp_step3}

Теперь лампой можно управлять и кнопкой, и командами --- оба способа работают одновременно,
потому что ни одна функция в \code{loop()} не останавливает программу.

\subsection{Прерывания}

Опрос в \code{loop()} работает, пока цикл крутится быстро. Если программа занята, короткое
нажатие можно пропустить. Альтернатива --- попросить процессор \emph{самого} вызвать функцию
в момент изменения уровня.

\begin{figure}[H]
\centering
\begin{tikzpicture}[x=1cm,y=1cm]
  \draw[thick,->] (0,0) -- (12,0) node[right,font=\small]{основная программа};
  \foreach \x in {0.3,1.3,...,11.3} \fill[accent!40] (\x,-0.15) rectangle (\x+0.8,0.15);
  \draw[very thick,esp,->] (4.7,1.3) node[above,font=\small]{фронт на GPIO} -- (4.7,0.2);
  \fill[esp!40] (4.8,-1.2) rectangle (6,-0.8) node[midway,font=\scriptsize]{ISR};
  \draw[->,esp] (4.7,-0.15) -- (4.8,-1.0);
  \draw[->,esp] (6,-1.0) -- (6.1,-0.15);
  \node[font=\scriptsize,align=center] at (8.5,-1.1) {обработчик: коротко!\\ установить флаг и выйти};
\end{tikzpicture}
\caption{Прерывание приостанавливает основную программу на время обработчика}
\end{figure}

\codecap{Подсчёт нажатий через прерывание}
\codeinput{l3_interrupt}

\begin{caution}
Внутри обработчика прерывания (ISR) нельзя вызывать \code{delay()}, \code{Serial.print()} и другие
медленные/блокирующие функции. Правило: изменить переменную --- и выйти. Всю работу делает основная программа.
\end{caution}

Для лампы опроса достаточно, поэтому в проекте оставляем \code{buttonPressed()}. К прерываниям
вернёмся в уроке~7, когда научимся будить из них задачи.

\subsection{Strapping-выводы}

Уровни на \pin{0}, \pin{3}, \pin{45}, \pin{46} считываются в момент сброса и определяют режим загрузки
и напряжение Flash. Не вешайте на них кнопки или нагрузку, которые могут изменить уровень при старте.
Именно поэтому кнопка лампы стоит на \pin{5}.

\begin{summary}
\begin{itemize}[leftmargin=*]
  \item Вход без подтяжки «висит»; используем \code{INPUT\_PULLUP} и кнопку на землю.
  \item Нажатие нужно ловить как \emph{событие} и фильтровать дребезг (\SIrange{20}{50}{ms}).
  \item Прерывание реагирует мгновенно, но обработчик должен быть минимальным.
\end{itemize}
\end{summary}

\begin{tasks}
\begin{enumerate}
  \item Используйте кнопку BOOT (\pin{0}) как вторую кнопку лампы после загрузки.
  \item Различайте короткое (< 0,5 с) и длинное (> 1 с) нажатие: длинное печатает в Serial время работы лампы.
  \item Замените кнопку сенсорным входом: \code{touchRead(2)} растёт при касании вывода \pin{2}.
\end{enumerate}
\end{tasks}
